\section{Conclusions}

The ARC-ADHD facial-feedback study is reframed as an incremental-information problem rather than an emotion-recognition problem. The module is justified only if within-child facial dynamics add robust out-of-sample information beyond task performance, context, self-report, and less intrusive telemetry.

Validation is staged: observational feature validation first, adaptive decision testing only after a prespecified incremental-value gate, and removal of the module if benefit is absent or fragile. This creates a falsifiable path for affective-computing research while preserving data minimization and child agency.
